\documentclass[11pt]{article}
\usepackage[margin=1in]{geometry}
\usepackage{amsmath,amssymb,amsthm}
\usepackage{mathtools}
\usepackage{booktabs}
\usepackage{array}
\usepackage{enumitem}
\usepackage{tikz}
\usepackage{hyperref}
\hypersetup{colorlinks=true,linkcolor=blue,citecolor=blue,urlcolor=blue}

\newtheorem{theorem}{Theorem}
\newtheorem{lemma}{Lemma}
\newtheorem{corollary}{Corollary}
\newtheorem{proposition}{Proposition}
\newtheorem{conjecture}{Conjecture}
\theoremstyle{definition}
\newtheorem{definition}{Definition}
\newtheorem{assumption}{Assumption}
\theoremstyle{remark}
\newtheorem{remark}{Remark}

\DeclareMathOperator{\dur}{dur}
\newcommand{\WA}{\mathsf{W_A}}
\newcommand{\WB}{\mathsf{W_B}}
\newcommand{\TT}{\mathsf{T}}
\newcommand{\NN}{\mathsf{N}}
\newcommand{\Term}{\mathrm{Term}}
\newcommand{\Z}{\mathbb{Z}}
\newcommand{\E}{\mathbb{E}}
\newcommand{\Tphi}{\mathcal{T}}

\title{On the Maximum Duration of a Deterministic\\ Card-Dueling Process\\[4pt]
\large A Combinatorial Analysis with Stochastic Generalizations}
\author{Ali Nehrani\\ Ground Inc.\\ \texttt{nehrani@gmail.com}}
\date{May 28, 2026 (revised, second revision)}

\begin{document}
\maketitle

\begin{abstract}
We study a deterministic two-player card-dueling process on two decks of $n$ cards each,
governed by a fixed pairwise outcome table. This process is a strategy-free generalization of the
classical card game War, in which a global look-up table replaces the linear ranking of cards and
a fourth ``no-interaction'' outcome is permitted. Our main result determines the exact extremal
duration: the maximum length of any terminating game is
\[
T(n) = \frac{n(n^2 + 2)}{3}.
\]
The proof combines a phase decomposition on a shrinking sequence of discrete tori $(\Z/m\Z)^2$
with the global perfect matching forced on the tie cells by termination; together these yield
a tight per-phase bound that sums to $T(n)$. The bound is attained by an explicit inductive
construction --- a Staircase Raster-Scan traversal --- whose induction is closed by a table-restriction
lemma. We then analyze stochastic generalizations, carefully separating two events
that an earlier treatment conflated: a random table admits some terminating ordering with
high probability, whereas for a fixed ordering termination is only a constant-probability event,
with $p^n \le \Pr(\Term_n) \le p$. Under the uniform model the expected number of \emph{distinct} duels
is $\Theta(n)$ (proved unconditionally), and the indicator-weighted total duration $\E[D_n\mathbf 1_{\Term_n}]$
is $O(n)$ \emph{under an explicit torus-coverage hypothesis}, which we reduce to a one-dimensional bound on the
longest monotone chain in the revealed set; this also yields a sub-quadratic $O(n^{3/2})$ bound under
a longest-increasing-subsequence-type estimate. The quadratic order $\Theta(n^2)$ lives in the sparse-tie
regime, where the tie density is $\Theta(1/n)$, and $\Theta(n^3)$ is approached only as the tie density tends to
$0$. Finally, a reproducible simulation (seeded, with bootstrap confidence intervals, $4\le n\le 18$, tie
density $p=\tfrac14$) shows that $\Pr(\Term_n)$ decays polynomially (least-squares fit $\approx 0.79\,n^{-2.35}$),
so that the conditional duration $E[D_n\mid \Term_n]$ grows \emph{super-linearly}, empirically as $\approx n^{1.4}$,
even though the indicator-weighted duration $\E[D_n\mathbf 1_{\Term_n}]$ is not merely bounded but
\emph{decreasing} (empirically $\approx n^{-0.95}$). This corrects an earlier claim that the conditional
duration is $\Theta(n)$; we emphasize that the fitted exponents rest on under one decade of $n$ and a single
value of $p$, and are reported as indicative rather than established.

\medskip
\noindent\textbf{Keywords:} deterministic deck dynamics; monovariants; perfect matchings; phase decomposition;
extremal combinatorics; random outcome tables; permutation statistics; card game of War.

\noindent\textbf{MSC 2020:} 05A05, 05C70, 60C05, 68Q25, 37B15.
\end{abstract}

\section{Introduction}

\subsection{Motivation}
Combinatorial processes that couple deterministic local update rules with global memory effects
recur throughout discrete mathematics and theoretical computer science. Sorting networks, the
abelian sandpile and chip-firing games, rotor-routing, cyclic cellular automata, and turn-based
card games all share a single signature: simple, local, memoryless transition rules generate global
behaviour whose duration, periodicity, or limiting shape is far from obvious and often requires an
extremal or potential-theoretic argument to pin down. The Card-Duel Process (CDP) studied in this
paper belongs squarely to this family, and it is most naturally read as a strategy-free generalization
of the card game War.

In War, two players each hold a face-down deck; each round both reveal a top card and the
higher-ranked card wins, with the loser (and sometimes both cards, on a tie) reinserted according
to fixed rules. War is famously deterministic once the initial shuffle is fixed: no player makes a
choice, yet the dynamics are intricate enough that even the basic question of whether the game must
terminate is subtle~\cite{LR}. The CDP abstracts War in two ways. First, the total order on card values
is replaced by an arbitrary outcome table $\phi : A \times B \to \{\WA, \WB, \TT, \NN\}$, decoupling the win relation
from any ranking and admitting genuinely non-transitive and asymmetric interactions. Second, a
fourth outcome $\NN$ (``no interaction'') is permitted, in which both cards cycle to the bottom without
a winner or loser; this is the move that allows the process to mark time on a torus without removing
mass. Removal happens only through ties, which makes the tie structure --- not the win structure ---
the carrier of all irreversibility in the process.

The driving question is purely extremal: across all outcome tables and all initial deck orderings,
how long can a terminating game last? Equivalently, we seek the worst-case running time of the
process as a function of the deck size $n$. For the concrete value $n = 40$ (one duel per minute) the
answer turns out to be exactly $356$ hours, a figure we recover as $T(40) = 21{,}360$.

\subsection{The model, informally}
Two players each hold a face-down deck of $n$ cards. Every step both reveal their top card; the
outcome of the duel is read off a fixed table $\phi$. On a win the winner stays on top and the loser goes
to the bottom of its own deck; on a no-interaction both cards go to the bottom of their decks; on a
tie both cards are permanently removed. No player has any choice at any point. The game ends
when both decks are empty. A precise statement is given in Section~\ref{sec:model}.

\subsection{Contributions}
\begin{itemize}[leftmargin=1.4em]
\item \textbf{An exact extremal duration (Theorem~\ref{thm:main}).} We prove that the maximum duration of any
terminating CDP on decks of size $n$ is exactly $T(n) = n(n^2 + 2)/3$. To our knowledge this is
the first sharp worst-case bound for a War-type process with an arbitrary outcome table.
\item \textbf{A phase decomposition with a tight per-phase bound (Section~\ref{sec:upper}).} We decompose any
game into $n$ phases delimited by ties and show that the $k$-th phase lives on a torus $(\Z/m_k\Z)^2$
with $m_k = n - k + 1$. The termination-forced perfect matching reserves $m_k - 1$ cells that the
phase provably cannot visit, giving $L_k \le m_k^2 - m_k + 1$; summation yields the upper bound.
\item \textbf{A matching explicit construction (Section~\ref{sec:lower}).} We exhibit, for every $n$, an outcome table
and initial ordering --- the Staircase Raster-Scan --- that attains $T(n)$. The construction is
recursive, and we close the induction with a table-restriction lemma (Lemma~\ref{lem:restrict}). The
general result is a theorem (Theorem~\ref{thm:tight}); we additionally checked the construction by
exhaustive simulation for $1 \le n \le 8$ as an independent confirmation.
\item \textbf{A corrected stochastic analysis (Section~\ref{sec:stoch}).} We separate the table-feasibility event $E_n$ (the
table admits some terminating ordering) from the fixed-ordering termination event $\Term_n$. We
prove $\Pr(E_n) \to 1$ exponentially fast, but $p^n \le \Pr(\Term_n) \le p$, so a fixed ordering terminates
only with constant probability. Conditioned on termination, the expected number of distinct
duels is $\Theta(n)$ (unconditional, Lemma~\ref{lem:fresh}), and the indicator-weighted total duration is $O(n)$
under an explicit torus-coverage hypothesis (Assumption~\ref{asm:cover}).
\item \textbf{A reduction of the coverage hypothesis (Section~\ref{sec:reduce}).} We prove (Theorem~\ref{thm:reduce}) that
Assumption~\ref{asm:cover} follows from a strictly weaker, one-dimensional condition --- that the revealed set
has no long monotone chain --- and that a longest-increasing-subsequence-type estimate already
yields a sub-quadratic $O(n^{3/2})$ bound (Proposition~\ref{prop:subquad}).
\item \textbf{Reproducible numerical findings (Section~\ref{sec:numerics}).} A seeded simulation with bootstrap
confidence intervals shows $\Pr(\Term_n) \approx 0.79\, n^{-2.35}$ (refuting the earlier sharp-termination
conjecture, now Conjecture~\ref{conj:poly}); $E[V \mid \Term_n] \approx 2.9\,n$ (linear, as predicted); a \emph{decreasing}
indicator-weighted duration $\E[D_n\mathbf 1_{\Term_n}] \approx 0.75\,n^{-0.95}$; and consequently a super-linear
conditional duration $E[D_n \mid \Term_n] \approx 0.95\,n^{1.40}$ --- correcting the earlier $\Theta(n)$ conditional
claim, but \emph{not} the $n^2$ that a constant-weighted-duration heuristic would predict
(Section~\ref{sec:numerics}).
\item \textbf{A map of the open landscape (Section~\ref{sec:open}).} We isolate a single coverage/mixing input ---
now in its reduced monotone-chain form --- on which several stochastic claims jointly depend,
and pose enumeration, termination-exponent, and generalization problems.
\end{itemize}

\subsection{Relation to previous versions}\label{sec:relation}
This is a revised version. We centralize all corrections here so that the technical sections read cleanly.

An early draft reported the expected duration under the uniform model as $\Theta(n^2)$. That error conflated
the number of \emph{matching} ties ($m$ per phase, a structural quantity) with the density of \emph{table}
ties ($\Theta(m^2)$ cells per phase under the uniform model). The first revision corrected this by (i)
distinguishing the table-feasibility event $E_n$ from the fixed-ordering termination event $\Term_n$;
(ii) replacing a flawed per-phase hitting-time argument by a global coupling that bounds only the number of
distinct cells; (iii) isolating the single torus-coverage hypothesis (Assumption~\ref{asm:cover}) on which the
total-duration order and the sparse-tie scaling both depend; (iv) reducing that hypothesis to a one-dimensional
monotone-chain bound (Theorem~\ref{thm:reduce}, Proposition~\ref{prop:subquad}); and (v) replacing the earlier
sharp-termination conjecture with the empirical claim $\Pr(\Term_n)\to 0$.

\emph{The present (second) revision} corrects an internal inconsistency in the numerical section of the first
revision. The first revision asserted that the indicator-weighted duration $\E[D_n\mathbf 1_{\Term_n}]$ ``stays
bounded'' and that the conditional duration $E[D_n\mid\Term_n]$ ``tracks $\approx n^2$.'' These two statements
are mutually inconsistent with the tabulated data and only both hold if $\E[D_n\mathbf 1_{\Term_n}]$ is a positive
constant. A reproducible simulation (Section~\ref{sec:numerics}) shows instead that $\E[D_n\mathbf 1_{\Term_n}]$ is
\emph{decreasing} (empirically $\approx n^{-0.95}$, hence $o(1)$, well inside the proven $O(n)$ envelope), and
consequently that $E[D_n\mid\Term_n]\approx n^{1.40}$ --- super-linear, but markedly sub-quadratic, not $n^2$.
This revision also (vi) reports bootstrap confidence intervals and a fixed random seed; (vii) re-interprets the
stale-traversal statistic $S$, whose growth and whose ratio $S/V$ to the fresh count both \emph{increase} with $n$,
so that the constant-$\kappa$ reading of Assumption~\ref{asm:cover} is not supported by the data and may fail
asymptotically; and (viii) flags that the fitted exponents rest on under one decade of $n$ and a single value of $p$.
The deterministic results (Sections~\ref{sec:phase}--\ref{sec:small}) are unchanged.

\subsection{Organization}
Section~\ref{sec:related} situates the CDP within the literature. Section~\ref{sec:model} fixes the model.
Section~\ref{sec:phase} develops the phase decomposition and the torus structure. Section~\ref{sec:upper} proves the
upper bound and Section~\ref{sec:lower} the matching lower bound. Section~\ref{sec:small} works small cases.
Section~\ref{sec:stoch} treats the stochastic generalizations. Section~\ref{sec:disc} discusses scope and limitations,
and Section~\ref{sec:open} lists open problems.

\section{Related Work}\label{sec:related}
The CDP intersects several established lines of research.

\paragraph{Card-game dynamics and the game of War.}
The closest relative of the CDP is the children's card game War, whose deterministic dynamics have received
careful attention. Lakshtanov and Roshchina~\cite{LR} showed that, contrary to a natural intuition, War can fail
to terminate for suitable reinsertion conventions, and characterized finiteness in terms of the reinsertion rule.
Our CDP replaces the rank order of War by an arbitrary outcome table and studies the orthogonal extremal
question --- the maximum finite duration. Bulgarian solitaire~\cite{Brandt} is another deterministic single-deck
card process whose orbit structure is governed by a monovariant.

\paragraph{Chip-firing, sandpiles, and rotor-routing.}
The abelian sandpile model and the chip-firing game~\cite{Dhar,Bjorner,LevinePeres} are paradigmatic examples
of deterministic local rules with strong global structure. The CDP shares the ``conserved-until-removed''
flavour. Rotor-routing~\cite{Holroyd} is closer in spirit: a deterministic rotor derandomizes a random walk yet
equidistributes, exactly the phenomenon we must invoke (Assumption~\ref{asm:cover}) when passing from the
deterministic walk on the torus to its stochastic counterpart. We caution, however, that we borrow only the
analogy: no quantitative coverage estimate is imported from the rotor-routing literature, and our
Assumption~\ref{asm:cover} is left as a hypothesis.

\paragraph{Cyclic cellular automata and annihilating particle systems.}
The $\NN$ and win moves act as cyclic rotations of each deck, placing the per-phase dynamics in the world of
cyclic update rules~\cite{Fisch,Bramson}; ties play the role of annihilation events. The CDP is a clean instance
in which the annihilation set is forced, by termination, to be a perfect matching.

\paragraph{Sorting and permutation dynamics.}
The win move ``loser to the bottom'' is the elementary operation underlying several classical sorting
analyses~\cite{Knuth}. The maximum-duration tables turn out to have non-tie entries arranged as a spiral-type
permutation pattern on each phase torus (Remark~\ref{rem:spiral}), loosely connecting the extremal construction
to nonattacking placements such as Costas arrays~\cite{Costas}. Equidistribution of the induced permutation walk
is governed by representation-theoretic tools in the spirit of~\cite{Diaconis,Stanley}.

\paragraph{Random matchings and termination thresholds.}
Our feasibility analysis (Proposition~\ref{prop:feas}) is a perfect-matching threshold for a random bipartite
graph whose edges are the tie cells; the relevant tools are standard~\cite{Bollobas,Janson,AlonSpencer,LP}. What
is new is the sharp distinction between a table possessing a matching (a graph property, high probability) and a
fixed ordering realizing a matching (a dynamical property, constant probability). For the empirical power-law
fitting in Section~\ref{sec:numerics} we follow the methodological cautions of Clauset, Shalizi, and
Newman~\cite{Clauset}.

\section{Model and Preliminaries}\label{sec:model}
Let $A = \{a_0, \dots, a_{n-1}\}$ and $B = \{b_0, \dots, b_{n-1}\}$ be disjoint sets of cards of equal size $n > 1$. A
configuration is a pair of ordered lists (decks) $(D_A, D_B)$ of the currently remaining cards, with $D_A[0]$
and $D_B[0]$ the respective top cards.

\begin{definition}[Outcome table]
An outcome table is a function $\phi : A \times B \to \{\WA, \WB, \TT, \NN\}$. Its tie-set is
$\Tphi(\phi) := \{(a, b) : \phi(a, b) = \TT\}$.
\end{definition}

\begin{definition}[Card-Duel Process]
Given an outcome table $\phi$ and an initial configuration $(D^0_A, D^0_B)$, the Card-Duel Process (CDP) evolves as
follows. At each step $t \ge 1$, let $a = D^{t-1}_A[0]$ and $b = D^{t-1}_B[0]$. According to $\phi(a, b)$:
\begin{itemize}[leftmargin=2.2em,nosep]
\item[$\WA$:] ($a$ wins) $a$ stays on top of $D_A$; $b$ moves to the bottom of $D_B$.
\item[$\WB$:] ($b$ wins) $b$ stays on top of $D_B$; $a$ moves to the bottom of $D_A$.
\item[$\TT$:] both $a$ and $b$ are permanently removed from their decks.
\item[$\NN$:] both $a$ and $b$ move to the bottoms of their decks.
\end{itemize}
The process terminates when both decks are empty.
\end{definition}

\begin{definition}[Game, duration]
A CDP is a \emph{game} if it terminates in finite time. We write $\mathcal G_n$ for the set of all games $(\phi, D^0)$
with $|A| = |B| = n$. For $G \in \mathcal G_n$, $\dur(G)$ denotes its duration: the total number of duels.
\end{definition}

\subsection{Basic observations}
\begin{lemma}[Only ties remove cards]\label{lem:ties}
In any CDP a card is removed if and only if it participates in a tie. Consequently, any game has exactly $n$ ties.
\end{lemma}
\begin{proof}
Outcomes $\WA, \WB, \NN$ return both involved cards to their decks; only $\TT$ removes cards. There are $2n$ cards and
each tie removes exactly one from each deck, so termination requires exactly $n$ ties.
\end{proof}

\begin{lemma}[Tie-matching]\label{lem:match}
For any game $G \in \mathcal G_n$, the pairs that tie during $G$ form a perfect matching between $A$ and $B$: there is a
bijection $\mu : A \to B$ with $\phi(a, \mu(a)) = \TT$ for all $a \in A$, and $(a, \mu(a))$ is the unique tie in which
$a$ participates.
\end{lemma}
\begin{proof}
By Lemma~\ref{lem:ties} each of the $2n$ cards is removed exactly once, namely by the unique tie in which it
participates. These $n$ ties therefore biject $A$ with $B$.
\end{proof}

The perfect matching $\mu$ is the single global object that drives the entire analysis: it is forced by termination,
it constrains every phase, and it is the structure that the stochastic sections randomize.

\section{Phase Decomposition and Torus Structure}\label{sec:phase}
\begin{definition}[Phase]
Phase $k$ ($1 \le k \le n$) is the maximal block of consecutive duels ending with the $k$-th tie of the game. Both
decks contain $m_k := n - k + 1$ cards throughout phase $k$.
\end{definition}

\begin{lemma}[Cyclic invariance]\label{lem:cyclic}
During any phase the relative cyclic order of cards within each deck is preserved.
\end{lemma}
\begin{proof}
Within a phase only $\WA, \WB, \NN$ occur, and each moves a top card directly to the bottom of its deck. Moving the
top element to the bottom is a pure cyclic rotation, so no card overtakes another and the cyclic word is invariant.
\end{proof}

\begin{definition}[Phase state and coordinate mapping]
At the start of phase $k$, label the cards currently in $A$ as $a_0, \dots, a_{m_k-1}$ from top to bottom, and likewise
$b_0, \dots, b_{m_k-1}$ in $B$. The phase state at a step $t$ within phase $k$ is the pair $(i, j) \in (\Z/m_k\Z)^2$
such that $a_i, b_j$ are the current top cards. By Lemma~\ref{lem:cyclic} the state space of phase $k$ is the discrete
torus $(\Z/m_k\Z)^2$.
\end{definition}

\begin{lemma}[Phase transition rule]\label{lem:trans}
Let $m = m_k$ and let $\phi_k$ denote $\phi$ restricted to the cards active in phase $k$. The state evolves from $(i, j)$ by
\begin{equation}\label{eq:trans}
(i, j) \longmapsto
\begin{cases}
(i,\; j + 1 \bmod m) & \phi_k(a_i, b_j) = \WA,\\
(i + 1 \bmod m,\; j) & \phi_k(a_i, b_j) = \WB,\\
(i + 1 \bmod m,\; j + 1 \bmod m) & \phi_k(a_i, b_j) = \NN,\\
\{\text{phase ends}\} & \phi_k(a_i, b_j) = \TT.
\end{cases}
\end{equation}
\end{lemma}
\begin{proof}
Direct transcription of Definition~2. If $\phi_k(a_i, b_j) = \WA$, card $a_i$ stays on top of $A$ while $b_j$ is sent to
the bottom of $B$; by Lemma~\ref{lem:cyclic} the card immediately below $b_j$ is $b_{j+1 \bmod m}$, giving $(i, j + 1)$.
The $\WB$ and $\NN$ cases are symmetric.
\end{proof}

The duration of phase $k$ is $L_k := \#\{\text{states visited in phase } k\}$, and $\dur(G) = \sum_{k=1}^n L_k$.

\begin{lemma}[No repeated states within a phase]\label{lem:saw}
For any game $G \in \mathcal G_n$, the trajectory within any phase visits each state at most once.
\end{lemma}
\begin{proof}
The transition rule~\eqref{eq:trans} is deterministic, so each non-tie state has a unique successor and predecessor
(each of the three non-tie maps is a bijection of $(\Z/m\Z)^2$). If a state were visited twice, the trajectory between
the two visits would be a directed cycle of non-tie states with no tie, hence repeat forever, contradicting finite
termination. For the final phase ($m_n = 1$) the torus is a single point, which must carry a tie for the game to end.
\end{proof}

\begin{remark}[Conceptual picture]
A game is thus a nested sequence of self-avoiding walks on tori of shrinking radius $m = n, n-1, \dots, 1$. Each walk
ends the instant it lands on a tie cell, which deletes one row and one column of the next torus. All the difficulty is
in how long a single walk can postpone hitting a tie.
\end{remark}

\section{The Upper Bound}\label{sec:upper}
\subsection{The phasewise bound}
\begin{theorem}[Phasewise bound]\label{thm:phasewise}
For every game $G \in \mathcal G_n$ and every phase $k$,
\begin{equation}\label{eq:phasewise}
L_k \le m_k^2 - m_k + 1, \qquad m_k = n - k + 1.
\end{equation}
\end{theorem}
\begin{proof}
Fix $G$ and $k$; write $m = m_k$ and let $S_k \subseteq (\Z/m\Z)^2$ be the set of states visited in phase $k$. By
Lemma~\ref{lem:saw}, $L_k = |S_k|$.

\emph{Future tie cells.} By Lemma~\ref{lem:match} the game has $n$ ties forming the perfect matching $\mu$. Exactly
$k - 1$ ties precede phase $k$, so $m = n - k + 1$ matching-ties remain: one ends phase $k$, and the other $m - 1$
occur in phases $k+1, \dots, n$. The $2(m - 1)$ cards involved in those future ties are all present at the start of
phase $k$, so each future tie corresponds to a cell $(i^*_\ell, j^*_\ell) \in (\Z/m\Z)^2$ with
$\phi_k(a_{i^*_\ell}, b_{j^*_\ell}) = \TT$. Because the matching is a bijection, distinct future ties use distinct
$a$-labels and distinct $b$-labels, hence distinct cells. Set
\begin{equation}
F_k := \{(i^*_2, j^*_2), \dots, (i^*_m, j^*_m)\}, \qquad |F_k| = m - 1.
\end{equation}

\emph{Future tie cells are never visited in phase $k$.} Suppose some $(i^*_\ell, j^*_\ell) \in F_k$ were visited during
phase $k$. By Lemma~\ref{lem:trans} the phase would end there, making $(a_{i^*_\ell}, b_{j^*_\ell})$ the $k$-th tie and
removing both cards in phase $k$. But these two cards are required, by the matching, for a tie in a strictly later
phase --- a contradiction. Hence $S_k \cap F_k = \emptyset$.

\emph{Counting.} Since $|(\Z/m\Z)^2| = m^2$ and $S_k \cap F_k = \emptyset$,
\[
L_k = |S_k| \le m^2 - |F_k| = m^2 - (m - 1) = m^2 - m + 1. \qedhere
\]
\end{proof}

\begin{remark}
The table $\phi$ may assign $\TT$ to cells beyond the $m$ matching-ties; such cells are also unvisitable in phase $k$,
so the bound~\eqref{eq:phasewise} is conservative. Using only the $m - 1$ guaranteed future-matching cells keeps the
argument purely structural; the extremal table places no surplus ties (Section~\ref{sec:lower}).
\end{remark}

\subsection{Summing over phases}
\begin{corollary}[Upper bound]\label{cor:upper}
For every $n \ge 1$ and every $G \in \mathcal G_n$,
\begin{equation}\label{eq:upper}
\dur(G) \le \sum_{k=1}^n \bigl(m_k^2 - m_k + 1\bigr) = \sum_{m=1}^n (m^2 - m + 1) = \frac{n(n^2 + 2)}{3}.
\end{equation}
\end{corollary}
\begin{proof}
Sum~\eqref{eq:phasewise} over $k$ and reindex by $m = n - k + 1$:
\[
\sum_{m=1}^n (m^2 - m + 1) = \frac{n(n+1)(2n+1)}{6} - \frac{n(n+1)}{2} + n
= \frac{n\bigl[(n+1)(2n+1) - 3(n+1) + 6\bigr]}{6}.
\]
Since $(n+1)(2n+1) - 3(n+1) + 6 = 2n^2 + 4$, this equals $\dfrac{n(2n^2 + 4)}{6} = \dfrac{n(n^2 + 2)}{3}$.
\end{proof}

\section{The Lower Bound: the Staircase Raster-Scan}\label{sec:lower}
We construct, for each $n$, an explicit table $\phi_n$ and initial configuration whose game attains $T(n)$. The
construction is recursive and the induction is closed by Lemma~\ref{lem:restrict}; the resulting tightness statement
(Theorem~\ref{thm:tight}) is a theorem valid for all $n$. As an independent check we verified by exhaustive
simulation that the construction attains $T(n)$ exactly for all $1 \le n \le 8$.

\subsection{Initial configuration and table}
Take the linear initial decks $D^{(n)}_A = [a_0, \dots, a_{n-1}]$ and $D^{(n)}_B = [b_0, \dots, b_{n-1}]$, so the phase-1
trajectory starts at $(0, 0)$. Place all $n$ tie cells on the anti-diagonal of the phase-1 torus:
\begin{equation}
\Tphi(\phi_n) = \bigl\{(i, j) \in (\Z/n\Z)^2 : i + j \equiv n - 1 \!\!\pmod n\bigr\}.
\end{equation}
The cell $(n-1, 0)$ (pairing $(a_{n-1}, b_0)$) is the active phase-ending tie, and the remaining anti-diagonal cells
form the forbidden set $F_1 = \{(i, n - 1 - i) : 0 \le i \le n - 2\}$.

The non-tie entries are assigned by a row-by-row staircase. For each row $i$ define
$c_N(i) = (n - 2 - i) \bmod n$ and $c_T(i) = (n - 1 - i) \bmod n$. Then:
\begin{itemize}[nosep]
\item Rows $i \in \{0, \dots, n - 2\}$: $\phi_n(a_i, b_{c_T(i)}) = \TT$, $\phi_n(a_i, b_{c_N(i)}) = \NN$, and
$\phi_n(a_i, b_j) = \WA$ for all other $j$.
\item Row $i = n - 1$: $\phi_n(a_{n-1}, b_0) = \TT$ and $\phi_n(a_{n-1}, b_j) = \WA$ for $j \in \{1, \dots, n - 1\}$.
\end{itemize}

\subsection{Canonical global-to-local restriction and dynamics}
Let phase $k$ have active size $m = n-k+1$, active cards
\[
A_m = \{a_0,\dots,a_{m-1}\},\qquad
B_m = \{b_{n-m},\dots,b_{n-1}\}.
\]
For $(a_i,b_j)\in A_m\times B_m$, define canonical local coordinates
\[
u=i,\qquad v=j-(n-m),\qquad 0\le u,v\le m-1.
\]
The global table restricts to the phase-local table
\[
\Phi_m(u,v)=
\begin{cases}
\TT & \text{if } u+v=m-1,\\
\NN & \text{if } u+v=m-2,\\
\WA & \text{otherwise}.
\end{cases}
\]
The local transition dynamics are
\[
\begin{aligned}
\Phi_m(u,v)=\WA &: (u,v)\mapsto\bigl(u,(v+1)\bmod m\bigr),\\
\Phi_m(u,v)=\NN &: (u,v)\mapsto\bigl((u+1)\bmod m,(v+1)\bmod m\bigr).
\end{aligned}
\]

\subsection{De-circularized row dynamics}
\begin{lemma}[Joint intermediate row-dynamics induction]\label{lem:joint-row}
Let $m>2$. Starting from $(0,0)$, for every row index $u\in\{0,\dots,m-2\}$:
\begin{enumerate}[label=(\roman*),nosep]
\item row $u$ is entered at $v_{\mathrm{start}}(u)=(m-u)\bmod m$;
\item row $u$ executes exactly $m-2$ consecutive $\WA$ steps and exits uniquely at $(u,m-2-u)$ via $\NN$;
\item that $\NN$ transition lands at $\bigl(u+1,(m-(u+1))\bmod m\bigr)$, the entry coordinate of row $u+1$.
\end{enumerate}
\end{lemma}
\begin{proof}
By induction on $u$.

For $u=0$, the entry is $(0,0)$ since $(m-0)\bmod m=0$. For $v\in\{0,\dots,m-3\}$, $0+v\le m-3$, so
the outcome is $\WA$ and $v$ increments by $1$. Thus exactly $m-2$ consecutive $\WA$ moves occur before
reaching $(0,m-2)$, where $0+(m-2)=m-2$, so the outcome is $\NN$. Its transition gives
\[
(0,m-2)\mapsto\bigl(1,(m-1)\bmod m\bigr),
\]
which is the row-1 entry coordinate.

Assume the statement for row $u-1$ with $1\le u\le m-2$. By the inductive handoff, row $u$ is entered at
$(u,(m-u)\bmod m)=(u,m-u)$. On columns $v\in\{m-u,\dots,m-1\}$ we have $u+v\ge m$, so outcomes are $\WA$.
After wrap-around, on columns $v\in\{0,\dots,m-3-u\}$ we have $u+v\le m-3$, so outcomes are also $\WA$.
These two segments have lengths $u$ and $m-2-u$, totaling $m-2$. The next column is $m-2-u$, where
$u+(m-2-u)=m-2$, so the unique exit is $\NN$ at $(u,m-2-u)$. Applying the $\NN$ transition yields
\[
\bigl((u+1)\bmod m,(m-(u+1))\bmod m\bigr),
\]
and because $u\le m-2$, this equals $\bigl(u+1,(m-(u+1))\bmod m\bigr)$.
\end{proof}

\begin{lemma}[Final-row traversal and termination]\label{lem:final-row}
On row $u=m-1$, the trajectory enters at column $v=1$, executes exactly $m-1$ consecutive $\WA$ steps,
and terminates at $(m-1,0)$ via $\TT$.
\end{lemma}
\begin{proof}
By Lemma~\ref{lem:joint-row}(iii) applied at $u=m-2$, entry to row $m-1$ is at
\[
v=(m-((m-2)+1))\bmod m=1.
\]
For $v\in\{1,\dots,m-1\}$, $(m-1)+v\ge m$, so outcomes are $\WA$, producing exactly $m-1$ steps and wrapping
to $(m-1,0)$. There, $(m-1)+0=m-1$, so the outcome is $\TT$ and the phase terminates.
\end{proof}

\begin{lemma}[Monotonic row development]\label{lem:row-mono}
From $(0,0)$, row coordinates are non-decreasing across the active phase timeline, progress through
$u=0,1,\dots,m-1$ in order, and terminate on row $m-1$ before any modular row wrap-around can occur.
\end{lemma}
\begin{proof}
By Lemma~\ref{lem:joint-row}, each intermediate row $u\in\{0,\dots,m-2\}$ has exactly one exit, an $\NN$ move to
row $u+1$ (without wrap since $u+1\le m-1$). By Lemma~\ref{lem:final-row}, row $m-1$ ends at a tie and stops.
Hence no transition after row $m-1$ occurs, so row wrap-around is impossible.
\end{proof}

\subsection{Trajectory injectivity and visited-row sets}
\begin{lemma}[Formal trajectory injectivity]\label{lem:traj-inj}
Let $\mathcal P_m=(S_0,S_1,\dots,S_{\tau_{\mathrm{term}}})$ be the local trajectory from $S_0=(0,0)$. If
$\tau\neq\tau'$, then $S_\tau\neq S_{\tau'}$.
\end{lemma}
\begin{proof}
Assume $\tau<\tau'$ and write $S_t=(u_t,v_t)$.

If $u_\tau\neq u_{\tau'}$, then by Lemma~\ref{lem:row-mono} we have $u_\tau<u_{\tau'}$, so states differ.

If $u_\tau=u_{\tau'}=u$, then both times lie within one row residency. In that residency, transitions are
$\WA$-moves until the unique row exit/terminal cell (Lemmas~\ref{lem:joint-row} and~\ref{lem:final-row}), so
$v\mapsto v+1\pmod m$ each step. The residency length is at most $m-1$ ($m-2$ on intermediate rows, $m-1$
on final row), so no column can repeat before leaving the row. Thus $v_\tau\neq v_{\tau'}$, hence
$S_\tau\neq S_{\tau'}$.
\end{proof}

\begin{definition}[Visited columns per row]\label{def:Vu}
For $u\in\{0,\dots,m-1\}$ define
\[
V_u=
\begin{cases}
\{0,1,\dots,m-2\}, & u=0,\\
\{m-u,\dots,m-1\}\cup\{0,\dots,m-2-u\}, & 1\le u\le m-2,\\
\{1,2,\dots,m-1,0\}, & u=m-1.
\end{cases}
\]
\end{definition}

\begin{lemma}[Exact cardinality of row sets]\label{lem:Vu-card}
\[
|V_u|=
\begin{cases}
m-1,& 0\le u\le m-2,\\
m,& u=m-1.
\end{cases}
\]
\end{lemma}
\begin{proof}
For $u=0$, $V_0=\{0,\dots,m-2\}$ has $m-1$ elements.

For $1\le u\le m-2$, the two intervals are disjoint and have sizes
$m-1-(m-u)+1=u$ and $(m-2-u)-0+1=m-1-u$, totaling $m-1$.

For $u=m-1$, the set contains all residues $0,1,\dots,m-1$, hence size $m$.
\end{proof}

\begin{theorem}[Complete global trajectory theorem]\label{thm:global-traj}
For any phase size $m>2$ initialized at $(0,0)$, the trajectory visits exactly the cells
$(u,v)$ with $v\in V_u$ once each, skips all non-terminal tie cells, and terminates at $(m-1,0)$ after
exactly $m^2-m+1$ visited states.
\end{theorem}
\begin{proof}
By Lemma~\ref{lem:joint-row}, row $0$ visits $V_0$ and exits at $(0,m-2)$; the omitted column is $m-1$, where
$0+(m-1)=m-1$ is a tie. For each $1\le u\le m-2$, Lemma~\ref{lem:joint-row} gives visited set
$V_u=\{m-u,\dots,m-1\}\cup\{0,\dots,m-2-u\}$ and unique omitted column $m-1-u$, where
$u+(m-1-u)=m-1$ is again tie. By Lemma~\ref{lem:final-row}, row $m-1$ visits
$V_{m-1}=\{1,\dots,m-1,0\}$ and terminates at $(m-1,0)$.

By Lemma~\ref{lem:traj-inj}, all visited states are distinct. Therefore
\[
L_1=\sum_{u=0}^{m-1}|V_u|
=(m-1)+\sum_{u=1}^{m-2}(m-1)+m
=(m-1)^2+m
=m^2-m+1.
\]
\end{proof}

\subsection{Base cases and inter-phase handoff}
\begin{lemma}[Base case $m=1$]\label{lem:base-m1}
The local phase has duration $1$ and terminates immediately; this equals $1^2-1+1$.
\end{lemma}
\begin{proof}
The only state is $(0,0)$ and $0+0=m-1$, so it is a tie and the phase ends at once.
\end{proof}

\begin{lemma}[Base case $m=2$]\label{lem:base-m2}
The local trajectory is
$(0,0)\xrightarrow{\NN}(1,1)\xrightarrow{\WA}(1,0)\xrightarrow{\TT}$,
so duration is $3=2^2-2+1$, skipping the non-terminal tie cell $(0,1)$.
\end{lemma}
\begin{proof}
Direct evaluation: $0+0=0=m-2$ gives $\NN$; then $1+1=2\notin\{0,1\}$ gives $\WA$; then $1+0=1=m-1$ gives $\TT$.
Hence exactly three visited states.
\end{proof}

\begin{lemma}[Deck-order permutation invariant and clean collapse]\label{lem:handoff}
Let phase $k$ have size $m$ with initial active sequences
\[
A_m=[a_0,\dots,a_{m-1}],\qquad B_m=[b_{n-m},\dots,b_{n-1}].
\]
At local state $(u_\tau,v_\tau)$, the physical decks are
\[
D_A^\tau=\bigl[a_{(u_\tau+\alpha)\bmod m}\mid \alpha=0,\dots,m-1\bigr],
\]
\[
D_B^\tau=\bigl[b_{n-m+(v_\tau+\beta)\bmod m}\mid \beta=0,\dots,m-1\bigr].
\]
At phase termination $(m-1,0)$, tie-removal deletes exactly $a_{m-1}$ and $b_{n-m}$, leaving
\[
[a_0,\dots,a_{m-2}],\qquad [b_{n-m+1},\dots,b_{n-1}],
\]
which is precisely the canonical start configuration of the next phase of size $m-1$.
\end{lemma}
\begin{proof}
Induct on local time $\tau$. The base case $\tau=0$ is immediate.

If the outcome is $\WA$, only deck $B$ rotates by one, so $v$ increments by one modulo $m$ and the displayed
formula remains true. If the outcome is $\NN$, both decks rotate by one, so both indices increment modulo $m$,
again preserving the formula.

By Theorem~\ref{thm:global-traj}, termination is at $(m-1,0)$, so
\[
D_A^{\mathrm{term}}=[a_{m-1},a_0,\dots,a_{m-2}],\qquad
D_B^{\mathrm{term}}=[b_{n-m},b_{n-m+1},\dots,b_{n-1}].
\]
Applying tie-removal deletes those top cards and yields exactly
$[a_0,\dots,a_{m-2}]$ and $[b_{n-m+1},\dots,b_{n-1}]$, the canonical phase-$(k+1)$ start.
\end{proof}

As an immediate phase-1 corollary of Theorem~\ref{thm:global-traj},
\begin{equation}\label{eq:L1}
L_1=n^2-n+1=m_1^2-m_1+1.
\end{equation}

\subsection{Closing the induction}
\begin{lemma}[Table restriction]\label{lem:restrict}
Let $\phi_n$ be the staircase table above. Identify the survivors of phase 1 via $a'_i = a_i$ ($0 \le i \le n-2$) and
$b'_j = b_{j+1}$ ($0 \le j \le n-2$). Then the restriction of $\phi_n$ to $\{a'_0, \dots, a'_{n-2}\} \times
\{b'_0, \dots, b'_{n-2}\}$ equals $\phi_{n-1}$ on the $(n-1)$-torus.
\end{lemma}
\begin{proof}
Write a surviving cell as $(i', j')$ with $i' = i \in \{0, \dots, n-2\}$ and $j' = j - 1$, $j \in \{1, \dots, n-1\}$.

\emph{Ties.} In $\phi_n$, ties satisfy $i + j \equiv n - 1 \pmod n$; for $i \le n - 2$, $j \ge 1$ this forces
$i + j = n - 1$. Thus $j' = n - 2 - i'$, which as $i'$ ranges over $\{0, \dots, n-2\}$ is exactly the anti-diagonal
$i' + j' \equiv n - 2 \pmod{n-1}$ of $\phi_{n-1}$; row $i' = n-2$ ties at $j' = 0$, matching the active tie of
$\phi_{n-1}$.

\emph{$\NN$ cells.} In $\phi_n$, row $i \le n - 2$ has its $\NN$ at $j = n - 2 - i$, which survives iff $i \le n - 3$;
for $i = n - 2$ the $\NN$ sits in the deleted column $0$. For $0 \le i' \le n - 3$ the surviving $\NN$ is at
$j' = n - 3 - i' = c^{(n-1)}_N(i')$ of $\phi_{n-1}$, and row $i' = n - 2$ carries no $\NN$, matching the final-row rule.

\emph{Remaining cells.} All other surviving cells are $\WA$ in $\phi_n$, and all non-tie non-$\NN$ cells of $\phi_{n-1}$
are $\WA$. The two tables agree on every cell.
\end{proof}

\begin{theorem}[Tightness]\label{thm:tight}
For every $n \ge 1$ the staircase game has duration exactly $T(n)$.
\end{theorem}
\begin{proof}
Induction on $n$. The base $n = 1$ is a single tie, $T(1) = 1$. For $n \ge 2$, phase 1 runs $L_1 = n^2 - n + 1$ duels
by~\eqref{eq:L1} and leaves the survivors in linear orders $D^{(n-1)}_A, D^{(n-1)}_B$. By Lemma~\ref{lem:restrict} the
survivors run under a table equal to $\phi_{n-1}$, so the remainder is the optimal $(n-1)$-game, lasting $T(n-1)$. Hence
\[
\dur = (n^2 - n + 1) + T(n - 1) = (n^2 - n + 1) + (n - 1)\frac{(n - 1)^2 + 2}{3} = \frac{n(n^2 + 2)}{3}. \qedhere
\]
\end{proof}

Corollary~\ref{cor:upper} and Theorem~\ref{thm:tight} together establish the main theorem.

\begin{theorem}[Main theorem]\label{thm:main}
For every $n \ge 1$, the maximum duration of a game is
\begin{equation}
\max_{G \in \mathcal G_n} \dur(G) = T(n) := \frac{n(n^2 + 2)}{3}.
\end{equation}
In particular $T(40) = 21{,}360$ duels; at one duel per minute this is $356$ hours.
\end{theorem}

\section{Small Cases}\label{sec:small}
\paragraph{$n = 1$.} One duel, a tie; $T(1) = 1$.

\paragraph{$n = 2$.} The table is $\phi_2 = \left(\begin{smallmatrix}\NN & \TT\\ \TT & \WA\end{smallmatrix}\right)$
(rows $a_0, a_1$; columns $b_0, b_1$). The trajectory is
$(0, 0) \xrightarrow{\NN} (1, 1) \xrightarrow{\WA} (1, 0) \xrightarrow{\TT}$, so $L_1 = 3$, then $L_2 = 1$; total
$T(2) = 4$.

\paragraph{$n = 3$.} The staircase table is
\[
\phi_3:\quad
\begin{array}{c|ccc}
 & b_0 & b_1 & b_2\\\hline
a_0 & \WA & \NN & \TT\\
a_1 & \NN & \TT & \WA\\
a_2 & \TT & \WA & \WA
\end{array}
\]
with tie-matching $\{(a_0, b_2), (a_1, b_1), (a_2, b_0)\}$. Phase 1 avoids $F_1 = \{(0, 2), (1, 1)\}$, visits
$3^2 - 3 + 1 = 7$ states, and ends at the active tie $(2, 0)$. The full 11-step trace is in Appendix~\ref{app:trace}.

\section{Stochastic Generalizations}\label{sec:stoch}
We now randomize the table. Throughout, the biased model $(\alpha, \beta, p, q)$ draws each entry of $\phi$
independently with $\Pr[\WA] = \alpha$, $\Pr[\WB] = \beta$, $\Pr[\TT] = p$, $\Pr[\NN] = q$, and
$\alpha + \beta + p + q = 1$; the uniform model is $\alpha = \beta = p = q = \tfrac14$.

Two distinct objects must be kept apart, as their conflation caused the error in the early draft (Section~\ref{sec:relation}).
The matching $\mu$ contributes exactly $m$ ties per phase (forced by termination), whereas the random table contains, in
expectation, $p\,m^2$ tie cells per phase. The trajectory ends a phase at the first tie cell it meets, so the relevant
rate is $p$ (the table-tie density), not $1/m$ (the matching density).

\subsection{The random duration}
Duration depends on both $\phi$ and the initial ordering $D^0$, so we fix the latter to the canonical linear orders
$D^0_A = [a_0, \dots, a_{n-1}]$, $D^0_B = [b_0, \dots, b_{n-1}]$ (any deterministic ordering behaves identically by
relabeling symmetry). Let $\Term_n$ be the event that the random game $(\phi, D^0)$ terminates, and let
$D_n = \dur(\phi, D^0)$ on $\Term_n$.

\subsection{Feasibility and termination probability}
\begin{definition}
$E_n$ is the set of tables $\phi$ that admit at least one terminating initial ordering.
\end{definition}

\begin{lemma}[Alignment sufficiency]\label{lem:align}
$\phi \in E_n$ iff $\Tphi(\phi)$ contains a perfect matching $A \to B$.
\end{lemma}
\begin{proof}
Necessity: a terminating ordering removes every card via a tie, so the realized ties form a perfect matching contained in
$\Tphi(\phi)$. Sufficiency: given a matching $\mu$, set $D^0_B = [\mu(a_0), \dots, \mu(a_{n-1})]$; then every duel is an
immediate tie and the game ends in $n$ steps.
\end{proof}

\begin{proposition}[Feasibility probability]\label{prop:feas}
Under the uniform model, $\Pr(E_n) \ge 1 - O\!\bigl(n(3/4)^n\bigr)$; hence $\Pr(E_n) \to 1$ exponentially fast.
\end{proposition}
\begin{proof}
By Lemma~\ref{lem:align}, $E_n^c$ is the event that the random bipartite graph $H = (A, B, E)$ with edges i.i.d.\ of
probability $p = \tfrac14$ has no perfect matching. By Hall's theorem this requires some $S \subseteq A$ with
$|S| = k$ and $|N(S)| \le k - 1$, giving $R \subseteq B$ of size $n - k + 1$ with no edges to $S$. With $q := \tfrac34$
and a union bound,
\[
\Pr(E_n^c) \le \sum_{k=1}^n \binom{n}{k}\binom{n}{k - 1} q^{k(n-k+1)}.
\]
The boundary terms $k = 1, n$ each equal $n q^n$, contributing $2 n q^n$. For the interior, $k(n - k + 1) = n + (k-1)(n-k)$
factors out $q^n$; the summand is symmetric under $k \mapsto n + 1 - k$, and using $\binom{n}{k}\binom{n}{k-1} \le n^{2k}$
one finds the interior contributes $o(n q^n)$. Hence $\Pr(E_n^c) = O\!\bigl(n(3/4)^n\bigr)$. (For constant $p$ the failure
of a perfect matching is dominated by isolated vertices, which gives the same $\Theta(n q^n)$ order.)
\end{proof}

A terminating run of $(\phi, D^0)$ realizes a tie-matching contained in $\Tphi(\phi)$, so $\Term_n \subseteq E_n$, and the
inclusion is strict. Consequently $\Pr(E_n) \to 1$ does \emph{not} imply $\Pr(\Term_n) \to 1$, and we bound
$\Pr(\Term_n)$ directly.

\begin{lemma}[Termination ceiling]\label{lem:ceiling}
For every $n \ge 1$, every biased model, and every fixed initial ordering, $\Pr(\Term_n) \le p$. In particular
$\Pr(\Term_n) \le \tfrac14$ under the uniform model.
\end{lemma}
\begin{proof}
On $\Term_n$ the realized ties form a perfect matching; let $(a^*, b^*)$ be the pair removed in the final phase $k = n$,
where $m_n = 1$. At the start of phase $n$ each deck holds the single card $a^*$ resp.\ $b^*$, and the state space is the
one point of $(\Z/1\Z)^2$; each of $\WA, \WB, \NN$ fixes this state, so the process terminates only if
$\phi(a^*, b^*) = \TT$.

Reveal table entries by deferred decisions as cells are first visited. The pair $(a^*, b^*)$ and the entire history through
phase $n - 1$ are measurable with respect to the cells visited in phases $1, \dots, n - 1$, and $(a^*, b^*)$ is not among
them (a prior visit revealing $\TT$ would have removed it earlier; a prior visit revealing a non-tie would freeze the cell
as a non-tie, precluding termination in phase $n$). Hence on $\Term_n$ the cell $(a^*, b^*)$ is fresh at phase $n$, and its
entry is independent of the history with $\Pr[\phi(a^*, b^*) = \TT] = p$. Therefore
$\Pr(\Term_n) \le p \cdot \Pr(\{(a^*, b^*)\text{ reached fresh}\}) \le p$.
\end{proof}

\begin{lemma}[Termination floor]\label{lem:floor}
For every $n \ge 1$ and every biased model, $\Pr(\Term_n) \ge p^n$.
\end{lemma}
\begin{proof}
If $\phi(a_i, b_i) = \TT$ for all $0 \le i \le n - 1$ (probability $p^n$), then under the canonical ordering every duel is
the diagonal tie $(a_i, b_i)$ and the game ends in $n$ steps. This event is contained in $\Term_n$.
\end{proof}

Thus $p^n \le \Pr(\Term_n) \le p$. The ceiling shows that termination of a fixed ordering is a constant-probability event,
never a high-probability one. An earlier version conjectured that the floor matches the ceiling in order,
$\Pr(\Term_n) = \Theta(p)$, bounded below by a positive constant independent of $n$. The supporting heuristic was that a
phase on the $m$-torus avoids a tie over $\Theta(m^2)$ fresh visits only with probability $(1 - p)^{\Theta(m^2)}$, so every
phase with $m$ above a constant ties with probability $\to 1$. Direct simulation contradicts this (Section~\ref{sec:numerics}):
under the uniform model $\Pr(\Term_n)$ decays polynomially. The flaw in the heuristic is that it counts only the per-phase
tie probability and ignores the $\Theta(n)$ independent opportunities, across the whole run, for the deterministic walk to
fall into a tie-free cycle (the true non-termination mode); these compound to drive $\Pr(\Term_n) \to 0$. We therefore
replace the earlier conjecture with the following.

\begin{conjecture}[Polynomial termination decay]\label{conj:poly}
Under the uniform model there is a constant $c > 0$ with $\Pr(\Term_n) = \Theta(n^{-c})$; in particular
$\Pr(\Term_n) \to 0$ polynomially. Our simulations suggest $c \approx 2.35$ for $p = \tfrac14$. We state this only for
$p = \tfrac14$, the single density we simulated; whether $\Pr(\Term_n) = \Theta(n^{-c(p)})$ for all fixed $p \in (0, 1)$,
and how $c$ depends on $p$, is open (Section~\ref{sec:open}, Q3). We give no theoretical derivation of the polynomial form,
which remains a numerically motivated guess.
\end{conjecture}

\subsection{Expected duration: distinct cells versus total duration}
An early draft bounded the per-phase length by a geometric first-tie time and summed to obtain
$\E[D_n \mathbf 1_{\Term_n}] \le n/p$. That argument is correct for the number of \emph{distinct} cells visited, but not
for the total duration, because a cell revealed (as a non-tie) in an earlier phase can be re-entered in a later phase,
adding a duel without carrying fresh tie-probability. We therefore separate the two quantities. Let $V$ denote the number
of distinct cells the walk visits over the whole game, and decompose
\begin{equation}\label{eq:decomp}
D_n = \sum_{k=1}^n L_k = V + S, \qquad S := \#\{\text{duels landing on an already-revealed cell}\}.
\end{equation}
Here $V = \sum_k F_k$ where $F_k$ is the number of cells first revealed in phase $k$, and $S = \sum_k A_k$ where $A_k$
counts the stale cells re-traversed in phase $k$. That $S$ can be positive is visible already in the deterministic $n = 3$
run of Appendix~\ref{app:trace}, where phase 2 re-traverses two cells revealed in phase 1.

\begin{lemma}[Fresh-cell coupling]\label{lem:fresh}
In the biased model with tie-probability $p \in (0, 1)$ and a fixed initial ordering, $\E[V \mathbf 1_{\Term_n}] \le n/p$.
\end{lemma}
\begin{proof}
Couple the table to two independent i.i.d.\ sequences: $(Y_t)_{t\ge1} \sim \mathrm{Bernoulli}(p)$ and $(U_t)_{t\ge1}$ on
$\{\WA, \WB, \NN\}$ with the conditional non-tie probabilities. Maintaining the set of revealed cells, on entering a
never-seen cell pop the next $(Y_t, U_t)$ and assign $\TT$ if $Y_t = 1$, else $U_t$; on re-entering a revealed cell reuse
its stored outcome. This samples $\phi$ with the correct marginals, and the tie indicator of the $t$-th fresh cell is
$Y_t$, independent of the walk's geometry through step $t - 1$.

Ties occur at fresh cells, so on $\Term_n$ the game ends at the $n$-th tie and $V = G_n$, where
$G_n := \inf\{t : \sum_{s \le t} Y_s = n\}$ is the $n$-th-success time. Thus $V \mathbf 1_{\Term_n} \le G_n$ pointwise,
and $\E[G_n] = n/p$ gives $\E[V \mathbf 1_{\Term_n}] \le n/p$.
\end{proof}

\begin{corollary}[Distinct-duel count]\label{cor:distinct}
Under the uniform model, $\E[V \mathbf 1_{\Term_n}] \le 4n$ (unconditionally). Conditioned on $\Term_n$ the expected
number of distinct duels lies between $n$ and $4n/\Pr(\Term_n)$; the upper bound is loose by a factor $\sim 1/\Pr(\Term_n)
\approx n^{2.35}$, so only the simulation (Section~\ref{sec:numerics}, $E[V \mid \Term_n] \approx 2.9\,n$) pins down the
constant. The rigorous content is the unconditional linear bound.
\end{corollary}

To pass from $V$ to the total duration $D_n$ we must control the stale re-traversals $S$. The trivial bound
$A_k \le |R_{<k}|$ only yields $S = O(nV) = O(n^2/p)$, and bounding the late phases crudely by $L_k \le m_k^2$ yields
$D_n = O(n^{3/2})$ for the uniform model; neither recovers the linear order. What is needed is:

\begin{assumption}[Bounded re-traversal / torus coverage]\label{asm:cover}
There is a constant $\kappa$ such that for every $n$ and every phase $k$, conditional on the cells revealed before phase
$k$, $\E[A_k \mid R_{<k}] \le \kappa\, \E[F_k \mid R_{<k}]$.
\end{assumption}

Assumption~\ref{asm:cover} holds whenever the phase walk equidistributes over its torus, so that the chance a step lands on
a stale cell is at most the stale density. In Section~\ref{sec:reduce} we reduce it to a strictly weaker one-dimensional
statement. \emph{We caution at the outset that the simulation of Section~\ref{sec:numerics} does not support the
constant-$\kappa$ form: both $\E[S \mid \Term_n]$ and the ratio $S/V$ \emph{increase} with $n$ over the accessible range,
so $\kappa$ is empirically not constant and may diverge.} All duration-order results below are therefore explicitly
conditional on Assumption~\ref{asm:cover}.

\begin{theorem}[Expected duration]\label{thm:dur}
Fix the canonical initial ordering and a biased model with tie-probability $p \in (0, 1)$.
\begin{enumerate}[label=(\roman*),nosep]
\item \emph{(Unconditional.)} The expected number of distinct cells visited satisfies $\E[V \mathbf 1_{\Term_n}] \le n/p$.
\item \emph{(Under Assumption~\ref{asm:cover}.)} The total expected duration satisfies
$\E[D_n \mathbf 1_{\Term_n}] \le (1 + \kappa)\dfrac{n}{p} = O\!\left(\dfrac{n}{p}\right)$. In particular, under the uniform
model $\E[D_n \mathbf 1_{\Term_n}] = O(n)$.
\end{enumerate}
\end{theorem}
\begin{proof}
(i) is Lemma~\ref{lem:fresh}. For (ii), summing Assumption~\ref{asm:cover} over phases gives
$\E[S \mathbf 1_{\Term_n}] \le \kappa\, \E[V \mathbf 1_{\Term_n}]$, whence by~\eqref{eq:decomp}
$\E[D_n \mathbf 1_{\Term_n}] \le (1 + \kappa)\,\E[V \mathbf 1_{\Term_n}] \le (1 + \kappa) n/p$.
\end{proof}

Combining with $D_n \ge n$ on $\Term_n$,
\[
n \Pr(\Term_n) \le \E[D_n \mathbf 1_{\Term_n}] \le (1 + \kappa)\,n/p \quad \text{(under Assumption~\ref{asm:cover})}.
\]
The conditional expectation $E[D_n \mid \Term_n] = \E[D_n \mathbf 1_{\Term_n}] / \Pr(\Term_n)$ is a different matter. An
earlier version asserted $E[D_n \mid \Term_n] = \Theta(n)$, but that conclusion silently assumed $\Pr(\Term_n) = \Theta(p)$,
which we now believe false (Conjecture~\ref{conj:poly}). With $\Pr(\Term_n) \to 0$ polynomially, the conditional duration is
super-linear; the empirical exponent (Section~\ref{sec:numerics}) is $\approx 1.40$, \emph{not} the $2$ that a constant
$\E[D_n\mathbf 1_{\Term_n}]$ would force, because the indicator-weighted duration is itself decreasing in $n$
(empirically $\approx n^{-0.95}$). Combined with the deterministic ceiling $T(n) = \Theta(n^3)$ we obtain, under
Assumption~\ref{asm:cover},
\begin{equation}\label{eq:interp}
\E[D_n \mathbf 1_{\Term_n}] = O\!\bigl(\min\{n/p, T(n)\}\bigr),
\end{equation}
as summarized in Table~\ref{tab:orders}.

\begin{table}[t]
\centering
\begin{tabular}{lll}
\toprule
tie density $p$ & order of $\E[D_n \mathbf 1_{\Term_n}]$ & comment\\
\midrule
$\Theta(1)$ (incl.\ uniform) & $O(n)$; empirically decreasing & each phase ties in $O(1)$ steps\\
$\Theta(1/n)$ & $O(n^2)$ (Conj.~\ref{conj:quad}) & the quadratic regime\\
$\Theta(1/n^2)$ & $O(n^3)$ & saturates the deterministic maximum\\
$p \to 1$ & $\to n$ & almost every duel is a tie\\
\bottomrule
\end{tabular}
\caption{Orders of $\E[D_n \mathbf 1_{\Term_n}]$ as the tie density $p$ varies, as upper bounds under
Assumption~\ref{asm:cover}; we claim no matching lower bounds. In the constant-$p$ row, simulation shows
$\E[D_n \mathbf 1_{\Term_n}]$ \emph{decreasing} ($\approx n^{-0.95}$), comfortably inside the $O(n)$ bound. The conditional
expectation $E[D_n \mid \Term_n]$ is obtained by dividing by $\Pr(\Term_n)$ and is therefore much larger when termination is
rare (Conjecture~\ref{conj:poly}). The distinct-cell bound $\E[V \mathbf 1_{\Term_n}] \le n/p$ is unconditional
(Lemma~\ref{lem:fresh}).}
\label{tab:orders}
\end{table}

\subsection{A reduction of the coverage hypothesis}\label{sec:reduce}
Assumption~\ref{asm:cover} is a two-dimensional equidistribution statement. We show it follows from a strictly weaker,
one-dimensional extremal condition: that the revealed set contains no long monotone chain.

\begin{lemma}[Geometric fresh budget]\label{lem:geom}
In the coupling of Lemma~\ref{lem:fresh}, the per-phase fresh counts $(F_k)$ are i.i.d.\ $\mathrm{Geometric}(p)$, so
$\E[F_k \mid R_{<k}] = 1/p$ for every phase whose torus $m_k$ exceeds a fixed constant; the remaining $O(1)$ small-torus
phases contribute $O(1)$ in total.
\end{lemma}
\begin{proof}
List the cells in order of first visit; their tie indicators are i.i.d.\ $\mathrm{Bernoulli}(p)$, independent of the walk's
geometry. A tie occurs only at a fresh cell, so the phase boundaries are the successes of this Bernoulli process and $F_k$
is the $k$-th inter-success gap. This requires fresh cells to remain available, which holds on every torus of size above a
constant since at most $O(k/p)$ cells have been revealed. For the $O(1)$ terminal phases of bounded torus size $m_k = O(1)$
the available fresh cells may be exhausted; there at least the surviving tie cell is fresh (revealing it earlier would have
removed the pair), so each such phase still ends, and the total contribution of these phases to both $V$ and $S$ is $O(1)$.
\end{proof}

Recall from Lemma~\ref{lem:trans} that consecutive visited cells differ by exactly one of $(0,1), (1,0), (1,1) \bmod m_k$;
call any such sequence a \emph{monotone chain}.

\begin{definition}[Stale set, longest monotone chain]
Let $\mathcal S_k$ be the set of card-pairs revealed before phase $k$ that survive into phase $k$, viewed on the phase-$k$
torus. Let $\Lambda(\mathcal S_k)$ be the length of the longest self-avoiding monotone chain contained in $\mathcal S_k$.
\end{definition}

\begin{lemma}[Run bound]\label{lem:run}
For every phase $k$, $A_k \le (F_k + 1)\,\Lambda(\mathcal S_k)$.
\end{lemma}
\begin{proof}
A maximal block of consecutive stale visits is, by the monotone-step rule and within-phase self-avoidance
(Lemma~\ref{lem:saw}), a self-avoiding monotone chain in $\mathcal S_k$, hence of length at most
$\Lambda(\mathcal S_k)$. Stale-runs are separated by fresh visits, of which there are $F_k$, so there are at most
$F_k + 1$ runs.
\end{proof}

\begin{theorem}[Coverage reduction]\label{thm:reduce}
Suppose there is a constant $C$ with $\E[\Lambda(\mathcal S_k) \mid R_{<k}] \le C$ for all $n$ and all $k$. Then
Assumption~\ref{asm:cover} holds with $\kappa = (1 + p)C$, and $\E[D_n \mathbf 1_{\Term_n}] \le \bigl(1 + (1 + p)C\bigr) n/p$.
\end{theorem}
\begin{proof}
$\Lambda(\mathcal S_k)$ is measurable w.r.t.\ $R_{<k}$, and by Lemma~\ref{lem:geom} $F_k$ is conditionally independent with
mean $1/p$. By Lemma~\ref{lem:run},
\[
\E[A_k \mid R_{<k}] \le \bigl(\E[F_k \mid R_{<k}] + 1\bigr)\E[\Lambda(\mathcal S_k) \mid R_{<k}]
\le \Bigl(\tfrac1p + 1\Bigr) C = (1 + p) C \cdot \tfrac1p = \kappa\, \E[F_k \mid R_{<k}].
\]
Summing over phases and adding $V$ gives the bound as in Theorem~\ref{thm:dur}(ii).
\end{proof}

The monotone-chain order is the coordinatewise partial order, so $\Lambda(\mathcal S_k)$ is a longest-chain
(Ulam--Hammersley) statistic; for a set behaving like $|\mathcal S_k|$ uniform points on the torus the longest chain has
order $2\sqrt{|\mathcal S_k|}$~\cite{AldousDiaconis,Romik}. This motivates a quantitative form.

\begin{proposition}[Sub-quadratic bound under a chain hypothesis]\label{prop:subquad}
If $\E[\Lambda(\mathcal S_k) \mid R_{<k}] \le c\sqrt{|\mathcal S_k|}$ for a constant $c$, then
$\E[D_n \mathbf 1_{\Term_n}] = O\!\bigl(n^{3/2} p^{-3/2}\bigr)$, which for the uniform model is $O(n^{3/2})$ --- strictly
better than the trivial $O(n^2)$.
\end{proposition}
\begin{proof}
Since $|\mathcal S_k| \le V$ and $\sqrt\cdot$ is concave,
$\E[\Lambda(\mathcal S_k)\mathbf 1_{\Term_n}] \le c\,\E[\sqrt V \mathbf 1_{\Term_n}] \le c\sqrt{\E[V \mathbf 1_{\Term_n}]}
\le c\sqrt{n/p}$. By Lemmas~\ref{lem:run} and~\ref{lem:geom}, $\E[A_k \mathbf 1_{\Term_n}] \le (1 + \tfrac1p) c \sqrt{n/p}$;
summing $n$ phases gives $\E[S \mathbf 1_{\Term_n}] = O(n^{3/2} p^{-3/2})$, and adding $\E[V \mathbf 1_{\Term_n}] \le n/p$
completes the bound.
\end{proof}

We do not prove the chain hypothesis: each earlier walk was monotone in a different modulus $m_j$, and re-reading it on the
phase-$k$ torus generically destroys monotonicity, so one expects $\Lambda(\mathcal S_k)$ to be small --- but establishing
this requires controlling the correlation between $\mathcal S_k$ and the phase-$k$ geometry, the residual content of Q2.
Notably, the simulation (Section~\ref{sec:numerics}) is consistent with $\Lambda(\mathcal S_k) = O(\sqrt{|\mathcal S_k|})$
(the longest stale-run grows slowly) yet shows $\E[S]$ growing super-linearly, which is exactly the regime in which
Proposition~\ref{prop:subquad} --- not the constant-$\kappa$ Assumption~\ref{asm:cover} --- is the relevant guarantee.

\subsection{Numerical findings}\label{sec:numerics}

\paragraph{Methodology and reproducibility.}
We simulated the uniform random-field CDP ($p = \tfrac14$) from the canonical ordering with a fixed pseudo-random seed
(\texttt{20260528}). Non-termination is detected \emph{exactly} rather than by a length cap: within a phase the relative
cyclic order is preserved (Lemma~\ref{lem:cyclic}), so the pair of top cards determines the phase state, and a repeat of
that pair since the last tie certifies a tie-free cycle (Lemma~\ref{lem:saw}), hence non-termination. This detection is
lossless; it is moreover \emph{independently} consistent with the deterministic ceiling, since by
Theorem~\ref{thm:main} no terminating run can exceed $T(n)$ steps, so a cap at $T(n)$ would yield identical statistics.
Each row of Table~\ref{tab:sim} pools the terminating runs from $2$--$6 \times 10^5$ trials (more for larger $n$ to keep
sample sizes up). Reported intervals are $95\%$ confidence intervals: Wilson for $\Pr(\Term_n)$, and $2000$-resample
bootstrap for the conditional means. The simulation code is included as supplementary material.

\begin{table}[t]
\centering
\footnotesize
\setlength{\tabcolsep}{4pt}
\begin{tabular}{rrrlllr}
\toprule
$n$ & trials & \#term & $\Pr(\Term_n)$ [95\% CI] & $E[D \mid \Term]$ [CI] & $E[V \mid \Term]$ [CI] & $E[S \mid \Term]$ [CI]\\
\midrule
4  & 200{,}000 & 5868 & 0.02934 [.0286,.0301] & 6.557 [6.50,6.61]  & 6.400 [6.36,6.44]  & 0.157 [.14,.17]\\
6  & 250{,}000 & 3025 & 0.01210 [.0117,.0125] & 11.693 [11.6,11.8] & 11.117 [11.0,11.2] & 0.576 [.53,.62]\\
8  & 300{,}000 & 1845 & 0.00615 [.0059,.0064] & 17.733 [17.5,18.0] & 16.548 [16.4,16.7] & 1.185 [1.1,1.3]\\
10 & 350{,}000 & 1313 & 0.00375 [.0036,.0040] & 24.381 [24.0,24.8] & 22.296 [22.0,22.6] & 2.085 [1.9,2.3]\\
12 & 400{,}000 & 881  & 0.00220 [.0021,.0024] & 31.722 [31.2,32.3] & 28.597 [28.2,29.0] & 3.125 [2.9,3.4]\\
14 & 450{,}000 & 750  & 0.00167 [.0016,.0018] & 37.844 [37.2,38.5] & 34.013 [33.6,34.5] & 3.831 [3.5,4.2]\\
16 & 500{,}000 & 589  & 0.00118 [.0011,.0013] & 45.902 [45.0,46.9] & 40.628 [40.0,41.2] & 5.273 [4.8,5.8]\\
18 & 600{,}000 & 521  & 0.00087 [.0008,.0010] & 54.015 [52.9,55.1] & 47.451 [46.7,48.1] & 6.564 [6.0,7.2]\\
\bottomrule
\end{tabular}
\caption{Uniform model ($p = \tfrac14$), conditioned on termination, fixed seed. $E[S \mid \Term]$ is the mean stale
re-traversal count ($S = D - V$). The decreasing indicator-weighted duration is $\E[D \mathbf 1_{\Term}] = E[D \mid \Term]
\cdot \Pr(\Term)$: $0.192, 0.142, 0.109, 0.092, 0.070, 0.063, 0.054, 0.047$ for $n = 4, \dots, 18$. The longest observed
stale-run over all runs at each $n$ is $5, 8, 9, 13, 12, 15, 17, 17$.}
\label{tab:sim}
\end{table}

\paragraph{Fitted scalings (log-log least squares over $4 \le n \le 18$).}
\begin{align*}
\Pr(\Term_n) &\approx 0.79\, n^{-2.35\,(\pm 0.03)}, &
E[D \mid \Term_n] &\approx 0.95\, n^{1.40\,(\pm 0.01)},\\
E[V \mid \Term_n] &\approx 2.94\, n - 6.4 \;\;(\text{affine; linear}), &
E[S \mid \Term_n] &\approx n^{2.43\,(\pm 0.11)},\\
\E[D_n \mathbf 1_{\Term_n}] &\approx 0.75\, n^{-0.95\,(\pm 0.03)}. & &
\end{align*}
The parenthesized values are regression standard errors on the exponent. We stress that $4 \le n \le 18$ spans under one
decade, so following~\cite{Clauset} these exponents should be read as indicative of order and sign, not as precisely
identified; in particular a slowly-decaying non-power form for $\Pr(\Term_n)$ cannot be excluded on this range, and only a
single tie density $p = \tfrac14$ was simulated.

Four observations emerge.
\begin{enumerate}[leftmargin=1.6em]
\item \textbf{Termination decays polynomially.} $\Pr(\Term_n) \approx 0.79\, n^{-2.35}$, inside the proven window
$p^n \le \Pr(\Term_n) \le p$ but \emph{not} bounded below; this is the basis for replacing the old sharp-termination
conjecture by Conjecture~\ref{conj:poly}. (The qualitative conclusion --- not bounded below --- is robust on this range;
the exponent $2.35$ is not.)

\item \textbf{Distinct duels are linear.} $E[V \mid \Term_n] \approx 2.9\, n$, comfortably inside the unconditional bound
$\E[V \mathbf 1_{\Term_n}] \le 4n$ of Lemma~\ref{lem:fresh} and confirming Corollary~\ref{cor:distinct}. This is the one
average-case order that is both proved (in the unconditional, indicator-weighted form) and observed.

\item \textbf{Re-traversals are \emph{not} negligible, and the constant-$\kappa$ assumption is unsupported.} The longest
\emph{single} stale-run grows slowly (reaching only $17$ at $n = 18$), consistent with
$\Lambda(\mathcal S_k) = O(\sqrt{|\mathcal S_k|})$ and hence with Proposition~\ref{prop:subquad}. But the \emph{total}
stale count $E[S \mid \Term_n]$ grows \emph{super-linearly} ($\approx n^{2.43}$), and the ratio $S/V$ rises
monotonically with $n$ (from $0.025$ at $n = 4$ to $0.138$ at $n = 18$, roughly $\propto n^{1.1}$). The strict
$\kappa = O(1)$ reading of Assumption~\ref{asm:cover} is therefore \emph{not} supported by the data and, extrapolated,
appears to fail: were $S/V$ to continue growing, stale traversals would eventually dominate and the total duration would
exceed linear order. (We caution that the largest-$n$ rows rest on $\sim 500$ terminating runs, so the precise exponent of
$E[S]$ is noisy; but the upward trend in $S/V$ is consistent across all rows.) The honest reading is that the linear total
duration of Theorem~\ref{thm:dur}(ii) has neither a proof nor clear empirical support at accessible sizes; the
$O(n^{3/2})$ guarantee of Proposition~\ref{prop:subquad} is the safer claim.

\item \textbf{Decreasing weighted duration; super-linear but sub-quadratic conditional duration.} The indicator-weighted
expectation $\E[D_n \mathbf 1_{\Term_n}]$ is not merely bounded but \emph{decreasing} ($\approx n^{-0.95}$, i.e.\ $o(1)$,
falling from $0.19$ to $0.05$ across the table), comfortably inside the $O(n)$ bound. The conditional duration
$E[D_n \mid \Term_n] = \E[D_n \mathbf 1_{\Term_n}] / \Pr(\Term_n)$ then grows like $n^{-0.95} / n^{-2.35} \approx n^{1.40}$,
matching the direct fit. This is super-linear --- correcting the earlier $\Theta(n)$ conditional claim --- but it is
\emph{markedly sub-quadratic}; the $\approx n^2$ that the first revision reported would follow only if
$\E[D_n \mathbf 1_{\Term_n}]$ were a positive constant, which the data refutes. The growth of the conditional duration is
driven almost entirely by the rarity of termination, not by long runs.
\end{enumerate}

\subsection{Where $\Theta(n^2)$ lives: the sparse-tie model}
By~\eqref{eq:interp} the quadratic order corresponds to $p = \Theta(1/n)$. A natural model realizing this is the
minimal-tie (sparse) model: the tie-set is exactly a uniformly random perfect matching $\mu$ (precisely $n$ ties), and
every non-tie cell is independent uniform on $\{\WA, \WB, \NN\}$. In phase $k$ the $m = n - k + 1$ survivors then expose
exactly $m$ tie cells on the $m \times m$ torus, an effective density $1/m$.

\begin{conjecture}[Quadratic scaling in the sparse-tie model]\label{conj:quad}
Under the minimal-tie model, conditioned on termination, $E[D_n] = \Theta(n^2)$, with $E[L_k] = \Theta(m)$ per phase.
\end{conjecture}
\begin{proof}[Heuristic]
A phase must locate one of $m$ targets among $m^2$ cells; if the walk equidistributes, the hitting time is $\Theta(m)$, and
$\sum_{m=1}^n \Theta(m) = \Theta(n^2)$. We state this as a conjecture because (i) the per-phase argument requires the same
coverage property as Assumption~\ref{asm:cover}, which we do not establish and which the uniform-model data calls into
question; and (ii) the phases are not independent. A rigorous proof must control both effects.
\end{proof}

\subsection{Low tie density}
\begin{remark}[$p \to 0$]
As $p \to 0$ the upper bound degrades and the duration is governed by the deterministic ceiling $T(n) = \Theta(n^3)$. We
caution against the stronger claim, made in an early draft, that the conditional duration converges to the exact extremal
value $T(n)$ as $p \to 0$: a generic sparse table produces a deterministic walk that typically enters a tie-free cycle
(hence non-termination) rather than the unique Hamiltonian staircase that realizes $T(n)$. Only the order $\Theta(n^3)$
(not the constant $1/3$) is plausibly recovered.
\end{remark}

\subsection{Combinatorial connections}
\begin{conjecture}[Count of extremal tables]\label{conj:count}
The number of extremal (maximum-duration) tables of size $n$ is $n! \cdot 2^{\Theta(n^2)}$.
\end{conjecture}
\noindent\emph{Status.} This is a heuristic. The factor $n!$ counts perfect tie-matchings; for each, the $n^2 - n$ non-tie
cells must produce a single self-avoiding phase-filling path, and a crude count of the directional degrees of freedom
suggests an exponent of order $n^2$. We have neither a lower nor an upper bound on the exponent constant, and it is
possible that the path constraints are restrictive enough to change the $2^{\Theta(n^2)}$ form; we record this as an open
problem (Q1), not a supported claim.

\begin{remark}[Spiral structure --- heuristic]\label{rem:spiral}
The staircase construction's non-tie entries trace a single Hamiltonian-type path that winds around each phase torus,
visually a spiral. We do not formalize a precise ``spiral permutation'' statistic here, and the resemblance to Costas
arrays~\cite{Costas} and to longest paths in de Bruijn graphs~\cite{Knuth} is offered only as informal motivation, not as a
theorem.
\end{remark}

\section{Discussion}\label{sec:disc}
\paragraph{What is proved, and what is assumed.}
The deterministic theory is complete and unconditional: Theorem~\ref{thm:main} pins down the exact extremal duration, with
a structural upper bound and an explicit, simulation-verified construction. The probabilistic theory is deliberately
stratified. \emph{Three statements are unconditional:} feasibility $\Pr(E_n) \to 1$ (Proposition~\ref{prop:feas}), the
termination window $p^n \le \Pr(\Term_n) \le p$ (Lemmas~\ref{lem:ceiling}--\ref{lem:floor}), and the linear bound on the
number of distinct duels (Lemma~\ref{lem:fresh}, Corollary~\ref{cor:distinct}). \emph{Everything concerning the total
duration order} --- the linear regime, the quadratic sparse-tie regime --- rests on the single torus-coverage input
(Assumption~\ref{asm:cover}). Section~\ref{sec:reduce} sharpens this: Assumption~\ref{asm:cover} reduces to bounding the
longest monotone chain (Theorem~\ref{thm:reduce}), and an $O(\sqrt\cdot)$ chain bound already gives the sub-quadratic
$O(n^{3/2})$ total. Simulation (Section~\ref{sec:numerics}) is the appropriate place to temper expectations: while the
longest single stale-run stays short (consistent with the $O(\sqrt\cdot)$ chain bound), the \emph{total} stale count and
the ratio $S/V$ both grow with $n$, so the constant-$\kappa$ Assumption~\ref{asm:cover} is not supported and the linear
total-duration claim should be regarded as conjectural; the $O(n^{3/2})$ bound is the conservative guarantee.

\paragraph{Why the distinction $E_n$ vs.\ $\Term_n$ matters.}
The early draft's error was conceptual: it implicitly assumed that a high probability of \emph{some} terminating ordering (a
graph property of the tie-set) transfers to a high probability that a \emph{fixed} ordering terminates (a dynamical
property of the walk). The ceiling $\Pr(\Term_n) \le p$ shows these are genuinely different: terminating from a fixed
shuffle is a constant-probability event because the final single-card phase must, against an independent fresh entry, land
on a tie. This is the same phenomenon that makes War's finiteness depend so delicately on the reinsertion rule~\cite{LR}.

\paragraph{Robustness.}
The phase decomposition uses only (i) that ties are the sole removal mechanism, (ii) that all other moves are cyclic
rotations, and (iii) finite termination. It is therefore robust to local changes in the win/no-interaction rules so long as
these preserve cyclic order within a phase, which makes the asymmetric-deck and higher-rank generalizations (Q5, Q6)
plausible targets.

\section{Open Problems}\label{sec:open}
\begin{description}[leftmargin=2.2em,style=nextline]
\item[Q1 (Enumeration).] Determine the exact number of extremal tables and the exponential generating function; settle
Conjecture~\ref{conj:count}, in particular the constant in the exponent.
\item[Q2 (Coverage, reduced form).] By Theorem~\ref{thm:reduce} it suffices to bound the longest monotone chain in the
revealed set: prove $\E[\Lambda(\mathcal S_k) \mid R_{<k}] = O(1)$ for Assumption~\ref{asm:cover}, $O(\sqrt{|\mathcal S_k|})$
for the sub-quadratic bound, or $o(|\mathcal S_k|)$ to beat $O(n^2)$. The obstruction is the correlation between
$\mathcal S_k$ and the phase-$k$ geometry. The simulation's growing $S/V$ ratio suggests the $O(1)$ form may be false;
clarifying which of these regimes holds is the central open question.
\item[Q3 (Termination exponent).] Prove $\Pr(\Term_n) = \Theta(n^{-c})$ for fixed $p$ (Conjecture~\ref{conj:poly}), identify
$c = c(p)$ (simulations suggest $c \approx 2.35$ at $p = \tfrac14$, on a single density and under one decade of $n$), and
characterize the tie-free-cycle mechanism that forces non-termination.
\item[Q4 (Average-case constants).] For the uniform model determine the constants in $E[V \mid \Term_n] = c_1 n + O(1)$
(simulations give $c_1 \approx 2.9$) and the exponent of $E[D_n \mid \Term_n]$ (simulations give $\approx 1.40$);
establish concentration.
\item[Q5 (Asymmetric decks).] Generalize to decks of sizes $n \ne m$ under modified one-sided removal rules; find
$\max \dur$ in terms of $n, m$.
\item[Q6 (Higher-rank duels).] Replace pairwise duels by $r$-wise interactions among $r \ge 3$ decks; find the analogue of
$T(n)$.
\end{description}

\section{Conclusion}
The maximum duration of the deterministic Card-Duel Process on two decks of $n$ cards is exactly $T(n) = n(n^2 + 2)/3$,
e.g.\ $21{,}360$ minutes ($356$ hours) for $n = 40$. The upper bound follows from a phase decomposition and the global
tie-matching forced by termination; the matching lower bound is realized by an explicit Staircase Raster-Scan construction
whose induction is closed by the table-restriction lemma. On the stochastic side, a random table admits some terminating
ordering with high probability ($\Pr(E_n) \to 1$), but for a fixed ordering termination is only a constant-probability
event, $p^n \le \Pr(\Term_n) \le p$ --- and in fact $\Pr(\Term_n)$ decays polynomially ($\approx 0.79\, n^{-2.35}$ in
simulation for $p = \tfrac14$), refuting an earlier constant-lower-bound conjecture. Conditioned on termination, the
expected number of distinct duels is $\Theta(n)$ under the uniform model (unconditionally), and the indicator-weighted
total duration is $O(n)$ under an explicit torus-coverage hypothesis --- which we reduce to a one-dimensional monotone-chain
bound and which already yields a sub-quadratic $O(n^{3/2})$ total under an Ulam--Hammersley-type estimate, and which the
simulation suggests may not hold in its strict constant-$\kappa$ form. Because $\Pr(\Term_n) \to 0$ while the
indicator-weighted duration is itself \emph{decreasing}, the conditional duration $E[D_n \mid \Term_n]$ is super-linear but
sub-quadratic, empirically $\approx n^{1.40}$ --- not the $\Theta(n)$ of an earlier draft, and not the $n^2$ that a
constant-weighted-duration heuristic would suggest. The most central remaining question is the reduced coverage estimate Q2
--- bounding the longest monotone chain in the revealed set --- on which several stochastic claims jointly depend.

\section*{Acknowledgments}
The author thanks the readers of previous drafts, whose questions exposed first the conflation of table-tie density with
matching density and then the inconsistency between the reported conditional-duration order and the simulation data, both of
which this revision corrects.

\begin{thebibliography}{99}
\bibitem{AldousDiaconis} D.~Aldous and P.~Diaconis. Hammersley's interacting particle process and longest increasing
subsequences. \emph{Probab. Theory Related Fields} 103(2):199--213, 1995.
\bibitem{AlonSpencer} N.~Alon and J.~H.~Spencer. \emph{The Probabilistic Method}, 4th ed. Wiley, 2016.
\bibitem{Bjorner} A.~Bj\"orner, L.~Lov\'asz, and P.~W.~Shor. Chip-firing games on graphs. \emph{European J. Combin.}
12(4):283--291, 1991.
\bibitem{Bollobas} B.~Bollob\'as. \emph{Random Graphs}, 2nd ed. Cambridge University Press, 2001.
\bibitem{Bramson} M.~Bramson and D.~Griffeath. Flux and fixation in cyclic particle systems. \emph{Ann. Probab.}
17(1):26--45, 1989.
\bibitem{Brandt} A.~Brandt. Cycles of partitions. \emph{Proc. Amer. Math. Soc.} 85(3):483--486, 1982.
\bibitem{Clauset} A.~Clauset, C.~R.~Shalizi, and M.~E.~J.~Newman. Power-law distributions in empirical data.
\emph{SIAM Review} 51(4):661--703, 2009.
\bibitem{Costas} J.~P.~Costas. A study of a class of detection waveforms having nearly ideal range--Doppler ambiguity
properties. \emph{Proc. IEEE} 72(8):996--1009, 1984.
\bibitem{Dhar} D.~Dhar. Self-organized critical state of sandpile automaton models. \emph{Phys. Rev. Lett.}
64(14):1613--1616, 1990.
\bibitem{Diaconis} P.~Diaconis. \emph{Group Representations in Probability and Statistics}. IMS Lecture Notes 11, 1988.
\bibitem{Fisch} R.~Fisch, J.~Gravner, and D.~Griffeath. Cyclic cellular automata in two dimensions. In \emph{Spatial
Stochastic Processes}, Progr. Probab. 19, pp.~171--185. Birkh\"auser, 1991.
\bibitem{Holroyd} A.~E.~Holroyd, L.~Levine, K.~M\'esz\'aros, Y.~Peres, J.~Propp, and D.~B.~Wilson. Chip-firing and
rotor-routing on directed graphs. In \emph{In and Out of Equilibrium 2}, Progr. Probab. 60, pp.~331--364. Birkh\"auser,
2008.
\bibitem{Janson} S.~Janson, T.~\L uczak, and A.~Ruci\'nski. \emph{Random Graphs}. Wiley-Interscience, 2000.
\bibitem{Knuth} D.~E.~Knuth. \emph{The Art of Computer Programming, Vol.~3: Sorting and Searching}, 2nd ed.
Addison-Wesley, 1998.
\bibitem{LR} E.~Lakshtanov and V.~Roshchina. On finiteness in the card game of War. \emph{Amer. Math. Monthly}
119(4):318--323, 2012.
\bibitem{LevinePeres} L.~Levine and Y.~Peres. Laplacian growth, sandpiles, and scaling limits. \emph{Bull. Amer. Math.
Soc.} 54(3):355--382, 2017.
\bibitem{LP} L.~Lov\'asz and M.~D.~Plummer. \emph{Matching Theory}. North-Holland, 1986.
\bibitem{Romik} D.~Romik. \emph{The Surprising Mathematics of Longest Increasing Subsequences}. Cambridge University
Press, 2015.
\bibitem{Stanley} R.~P.~Stanley. \emph{Enumerative Combinatorics, Vol.~1}, 2nd ed. Cambridge University Press, 2012.
\end{thebibliography}

\appendix
\section{Numerical Values of $T(n)$}\label{app:tn}
From $T(n) = n(n^2 + 2)/3$ (one duel per minute):
\begin{center}
\begin{tabular}{rrrr}
\toprule
$n$ & $T(n)$ (duels) & minutes & hours\\
\midrule
1 & 1 & 1 & 0.017\\
2 & 4 & 4 & 0.067\\
3 & 11 & 11 & 0.183\\
4 & 24 & 24 & 0.400\\
5 & 45 & 45 & 0.750\\
10 & 340 & 340 & 5.667\\
20 & 2{,}680 & 2{,}680 & 44.67\\
30 & 9{,}010 & 9{,}010 & 150.2\\
40 & 21{,}360 & 21{,}360 & 356.0\\
50 & 41{,}700 & 41{,}700 & 695.0\\
100 & 333{,}400 & 333{,}400 & 5{,}557.0\\
\bottomrule
\end{tabular}
\end{center}

\section{Full Trace for the $n = 3$ Staircase Game}\label{app:trace}
Initial $D^0_A = [a_0, a_1, a_2]$, $D^0_B = [b_0, b_1, b_2]$, table $\phi_3$ as in Section~\ref{sec:small}.
\begin{center}
\footnotesize
\begin{tabular}{rlllll}
\toprule
Duel & $D_A$ & $D_B$ & Top & Outcome & Action\\
\midrule
\multicolumn{6}{l}{\emph{Phase 1 ($m = 3$)}}\\
1 & $[a_0,a_1,a_2]$ & $[b_0,b_1,b_2]$ & $(a_0,b_0)$ & $\WA$ & $b_0 \to$ bottom\\
2 & $[a_0,a_1,a_2]$ & $[b_1,b_2,b_0]$ & $(a_0,b_1)$ & $\NN$ & both $\to$ bottom\\
3 & $[a_1,a_2,a_0]$ & $[b_2,b_0,b_1]$ & $(a_1,b_2)$ & $\WA$ & $b_2 \to$ bottom\\
4 & $[a_1,a_2,a_0]$ & $[b_0,b_1,b_2]$ & $(a_1,b_0)$ & $\NN$ & both $\to$ bottom\\
5 & $[a_2,a_0,a_1]$ & $[b_1,b_2,b_0]$ & $(a_2,b_1)$ & $\WA$ & $b_1 \to$ bottom\\
6 & $[a_2,a_0,a_1]$ & $[b_2,b_0,b_1]$ & $(a_2,b_2)$ & $\WA$ & $b_2 \to$ bottom\\
7 & $[a_2,a_0,a_1]$ & $[b_0,b_1,b_2]$ & $(a_2,b_0)$ & $\TT$ & remove $a_2,b_0$\\
\multicolumn{6}{l}{\emph{Phase 2 ($m = 2$)}}\\
8 & $[a_0,a_1]$ & $[b_1,b_2]$ & $(a_0,b_1)$ & $\NN$ & both $\to$ bottom\\
9 & $[a_1,a_0]$ & $[b_2,b_1]$ & $(a_1,b_2)$ & $\WA$ & $b_2 \to$ bottom\\
10 & $[a_1,a_0]$ & $[b_1,b_2]$ & $(a_1,b_1)$ & $\TT$ & remove $a_1,b_1$\\
\multicolumn{6}{l}{\emph{Phase 3 ($m = 1$)}}\\
11 & $[a_0]$ & $[b_2]$ & $(a_0,b_2)$ & $\TT$ & game over\\
\bottomrule
\end{tabular}
\end{center}
The run lasts $7 + 3 + 1 = 11 = T(3)$ duels, visiting no coordinate twice within a phase. Note, however, that cells
$(0, 1)$ and $(1, 2)$ are visited in both phase 1 and phase 2: these cross-phase re-traversals are the stale visits $S$
of~\eqref{eq:decomp}, present even in the deterministic optimum.

\end{document}